\section{Формула Остроградского"--~Гаусса}

\begin{defn}
Область $G \subset \bbR^3_{xyz}$ называется \textit{канонической\rindex{область!каноническая в $\bbR^3$}}, если она одновременно задаётся как неравенствами
$$
\psi(x, y) < z < \varphi(x, y), \quad (x, y) \in G_1
$$
($G_1$ "--- измеримая область в $\bbR^2_{xy}$), где функции $\psi$ и $\varphi$ непрерывны в $\overline{G}_1$, так и неравенствами
$$
\xi(y, z) < x < \eta(y, z), \quad (y, z) \in G_2
$$
($G_2$ "--- измеримая область в $\bbR^2_{yz}$), где функции $\xi$ и $\eta$ непрерывны в $\overline{G}_2$, а также неравенствами
$$
\alpha(z, x) < y < \beta(z, x), \quad (z, x) \in G_3
$$
($G_3$ "--- измеримая область в $\bbR^2_{zx}$), где функции $\alpha$ и $\beta$ непрерывны в $\overline{G}_3$.

Замыкание канонической области является простым множеством одновременно относительно всех трёх координатных осей.
\end{defn}

Формула Остроградского"--~Гаусса связывает значение поверхностного интеграла второго рода по границе области в $\bbR^3$ со значением некоторого тройного интеграла по всей области.

\begin{thm}[формула Остроградского"--~Гаусса]
Пусть $G$ "--- каноническая область в пространстве $\bbR^3$, векторное поле $\vec{a}(x, y, z) = (P, Q, R)^T$ непрерывно дифференцируемо в $\overline{G}$. Тогда
$$
\iiint\limits_G \div \vec{a}\,dx\,dy\,dz = \iint\limits_{\partial G} (\vec{a}, d\vec{S}), \quad \text{т.е.}
$$
$$
\iiint\limits_G \left( \pd{P}{x} + \pd{Q}{y} + \pd{R}{z} \right) dx\,dy\,dz = \iint\limits_{\partial G} \bigl( P\,dy\,dz + Q\,dz\,dx + R\,dx\,dy \bigr)
$$
(граница области ориентирована внешней нормалью).
\end{thm}

\begin{notion}
Доказательство теоремы проходит для более общего случая, когда $P, Q, R$ непрерывны в $\overline{G}$, $\pd{P}{x}$, $\pd{Q}{y}$, $\pd{R}{z}$ непрерывны и ограничены в $G$, причём существуют конечные пределы
\begin{equation*}
\begin{aligned}
&\lim_{z \to \varphi(x,y)-0} \pd{R}{z}, \quad \lim_{z \to \psi(x,y)+0} \pd{R}{z}, \quad (x, y) \in \overline{G}_1, \\
&\lim_{x \to \eta(y,z)-0} \pd{P}{x}, \quad \lim_{x \to \xi(y,z)+0} \pd{P}{x}, \quad (y, z) \in \overline{G}_2, \\
&\lim_{y \to \beta(z,x)-0} \pd{Q}{y}, \quad \lim_{y \to \alpha(z,x)+0} \pd{Q}{y}, \quad (z, x) \in \overline{G}_3
\end{aligned}
\end{equation*}
(обозначения как в определении 21.4).
\end{notion}

\begin{proof}
По теореме о сведении кратного интеграла к повторному, применённой к $\overline{G}$ как к простому множеству относительно оси $z$,
$$
\iiint\limits_{\overline{G}} \pd{R}{z}\,dx\,dy\,dz = \iint\limits_{\overline{G}_1} dx\,dy \int\limits_{\psi(x,y)}^{\varphi(x,y)} \pd{R(x, y, z)}{z}\,dz.
$$
Тройной интеграл существует по теореме 19.5 (в более общем случае по теореме 19.9). Так как граница области $G$ состоит из графиков непрерывных функций $z = \varphi(x, y)$ и $z = \psi(x, y)$ на измеримом компакте $\overline{G}_1$, а также части цилиндрической поверхности, то $\mu \partial G = 0$. Тройной интеграл по множеству $G$ такой же, как и по множеству $\overline{G}$.

При фиксированных $(x, y) \in G_1$ функция $R(x, y, z)$ непрерывна по $z$ на отрезке $[\psi(x, y); \varphi(x, y)]$, непрерывно дифференцируема на интервале $(\psi(x, y); \varphi(x, y))$ и $\pd{R}{z}$ имеет конечные пределы в концах интервала. С учётом теоремы 4.16 эта функция непрерывно дифференцируема на $[\psi(x, y); \varphi(x, y)]$, и по формуле Ньютона"--~Лейбница
$$
\int\limits_{\psi(x,y)}^{\varphi(x,y)} \pd{R(x, y, z)}{z}\,dz = R(x, y, \varphi(x, y)) - R(x, y, \psi(x, y)).
$$

Применяя сведение тройного интеграла к повторному, имеем
\begin{multline*}
\iiint\limits_G \pd{R}{z}\,dx\,dy\,dz = \iint\limits_{\overline{G}_1} dx\,dy \int\limits_{\psi(x,y)}^{\varphi(x,y)} \pd{R(x, y, z)}{z}\,dz = \\
= \iint\limits_{\overline{G}_1} R(x, y, \varphi(x, y))\,dx\,dy - \iint\limits_{\overline{G}_1} R(x, y, \psi(x, y))\,dx\,dy = \\
= \iint\limits_{\Gamma_\varphi \uparrow} R(x, y, z)\,dx\,dy - \iint\limits_{\Gamma_\psi \uparrow} R(x, y, z)\,dx\,dy,
\end{multline*}
где последние два поверхностных интеграла второго рода берутся по верхним сторонам графиков соответствующих непрерывных функций на компакте $\overline{G}_1$ (см. определение 20.17). Внешняя сторона $\partial G$ соответствует верхней стороне графика функции $\varphi$ и нижней стороне графика функции $\psi$, поэтому в последнее выражение интегралы по внешним сторонам соответствующих кусков $\partial G$ войдут со знаком <<$+$>> (см. рис. 21.1).

Остался нерассмотренным интеграл $\iint\limits_S R(x, y, z)\,dx\,dy$, где $S$ "--- часть цилиндрической поверхности, входящая в $\partial G$ (см. рис. 21.1). Но для этого интеграла соответствующая вектор-функция $\vec{a} = (0, 0, R)^T$, а вектор $\vec{\nu}$ нормали к $S$ ортогонален оси $z$, поэтому $(\vec{a}, \vec{\nu}) = 0$, и
$$
\iint\limits_S R\,dx\,dy = \iint\limits_S (\vec{a}, \vec{\nu})\,dS = 0.
$$

Окончательно имеем
$$
\iiint\limits_G \pd{R}{z}\,dx\,dy\,dz = \iint\limits_{\partial G} R(x, y, z)\,dx\,dy,
$$
где поверхностный интеграл второго рода берётся по внешней стороне $\partial G$. Аналогично, рассматривая $\overline{G}$ как простое множество относительно осей $x$ и $y$, получим
$$
\iiint\limits_G \pd{P}{x}\,dx\,dy\,dz = \iint\limits_{\partial G} P(x, y, z)\,dy\,dz,
$$
$$
\iiint\limits_G \pd{Q}{y}\,dx\,dy\,dz = \iint\limits_{\partial G} Q(x, y, z)\,dz\,dx.
$$
Складывая полученные три равенства, придём к нужной формуле.
\end{proof}

Поверхностный интеграл второго рода в теории поля называют \textit{потоком\rindex{поток векторного поля}} векторного поля через выбранную сторону поверхности. Таким образом, поток непрерывно дифференцируемого векторного поля через внешнюю сторону границы канонической области в $\bbR^3$ равен тройному интегралу от дивергенции поля по всей области.

Имеет место следующее обобщение этой теоремы.

\begin{thm}[21.1$'$]
Пусть $G$ "--- ограниченное открытое множество в $\bbR^3$, граница $\partial G$ которого состоит из конечного числа КГП. При этом $G$ разбивается на конечное число канонических областей $G_1, \ldots, G_N$. Это значит, что $\overline{G}_1, \ldots, \overline{G}_N$ образуют разбиение $\overline{G}$; разные $G_i$ не имеют общих точек, а $\overline{G}_i$ могут иметь лишь общие участки границ. Тогда
$$
\iiint\limits_G \div \vec{a}\,dx\,dy\,dz = \iint\limits_{\partial G} (\vec{a}, d\vec{S}),
$$
если векторное поле $\vec{a}$ непрерывно дифференцируемо в $\overline{G}$ (граница $G$ ориентирована внешней нормалью).
\end{thm}


\section{Соленоидальные векторные поля}

\begin{defn}
Векторное поле $\vec{a}$ называется \textit{соленоидальным\rindex{векторное поле!соленоидальное}} в области $G \subset \bbR^3$, если поток этого поля через любую замкнутую КГП $S \subset G$ равен нулю.
\end{defn}

\begin{thm}
Если непрерывно дифференцируемое векторное поле $\vec{a}$ соленоидально в области $G \subset \bbR^3$, то $\div \vec{a} = 0$ в области $G$. Обратно, если в объёмно односвязной области $G \subset \bbR^3$ для непрерывно дифференцируемого векторного поля $\vec{a}$ выполняется равенство $\div \vec{a} = 0$, то поле $\vec{a}$ соленоидально в $G$.
\end{thm}

\begin{proof}
Если точка $M_0 \in G$, то по теореме 21.2 о геометрическом смысле дивергенции
$$
\div \vec{a}(M_0) = \lim_{k \to \infty} \frac{\iint\limits_{\partial G_k} (\vec{a}, d\vec{S})}{\mu G_k},
$$
где в качестве $G_k$ можно взять последовательность окрестностей точки $M_0$, радиусы которых стремятся к нулю. Так как все интегралы в этом равенстве равны нулю при $G_k \subset G$ (т.е. при $k \geqslant k_0$), то $\div \vec{a}(M_0) = 0$, а $M_0$ "--- произвольная точка $G$. Первая часть теоремы доказана.

Обратно, пусть $\div \vec{a} = 0$ в объёмно односвязной области $G$, а $S$ "--- замкнутая КГП в $G$. В силу объёмной односвязности $G$ поверхность $S$ является границей ограниченного открытого множества $D \subset G$. По формуле Остроградского"--~Гаусса (мы приняли без доказательства, что формула Остроградского"--~Гаусса справедлива в этом общем случае)
$$
\iint\limits_S (\vec{a}, d\vec{S}) = \iiint\limits_D \div \vec{a}\,dx\,dy\,dz = 0
$$
(поверхность $S$ ориентирована внешней нормалью). Так как $\iint\limits_S (\vec{a}, d\vec{S}) = 0$ по любой замкнутой КГП в $G$, то поле $\vec{a}$ соленоидально в области $G$.
\end{proof}

\begin{notion}
Условие объёмной односвязности области $G$ во второй части теоремы существенно.
\end{notion}